\section{Programm: Fibonacci}
\label{sec:c:fibonacci}

Die Fibonacci-Zahlen\footnote{siehe \cite{wikipedia:fibonacci}} sind eine Folge von Zahlen, bei der jede Zahl die Summe der beiden vorhergehenden ist. Die Folge beginnt mit den Zahlen 0 und 1. Die ersten paar Fibonacci-Zahlen sind:
\begin{align*}
  F(0) &= 0, \\
  F(1) &= 1, \\
  F(2) &= 1, \\
  F(3) &= 2, \\
  F(4) &= 3, \\
  F(5) &= 5, \\
  F(6) &= 8, 
\end{align*}
und so weiter. Die Fibonacci-Zahlen haben viele interessante Eigenschaften und Anwendungen in der Mathematik, der Natur und der Informatik. Sie tauchen in der Natur auf, zum Beispiel in der Anordnung von Blättern an einem Stängel, in der Verzweigung von Bäumen, in der Anordnung von Blütenblättern und in der Struktur von Tannenzapfen. 

Die Fibonacci-Zahlen werden folgendermaßen berechnet: Die ersten beiden Zahlen der Folge sind 0 und 1. Jede weitere Zahl ist die Summe der beiden vorhergehenden Zahlen. Hier ist ein Beispielprogramm, das die ersten 16 Fibonacci-Zahlen berechnet und ausgibt (die 0 als Startwert lassen wir weg).

\cppinput{c/03_fibonacci.c}

\paragraph{Erklärung:}

\begin{description}
  \item[\texttt{int a = 0;}] \hfill\\
    Wir initialisieren die erste Fibonacci-Zahl mit 0.
  \item[\texttt{int b = 1;}] \hfill\\
    Wir initialisieren die zweite Fibonacci-Zahl mit 1.
  \item[\texttt{for(int i = 1; i <= 16; i = i+1) \{ … \} }] \hfill\\
    Wir starten eine Schleife, die 16 Mal durchläuft, um die ersten 16 Fibonacci-Zahlen zu berechnen. Wisst Ihr übrigens, dass es in diesem Buch auch ein ausführliches Nachschlagewerk gibt? Dort könnt Ihr nachlesen, wie Schleifen genau funktionieren — schaut mal in Abschnitt~\ref{sec:referenz:c:schleifen} nach.
  \item[\texttt{int c = a + b;}] \hfill\\
    Wir berechnen die nächste Fibonacci-Zahl, indem wir die beiden vorhergehenden Zahlen addieren.
  \item[\texttt{a = b;}] \hfill\\
    Wir verschieben die Werte: Die aktuelle zweite Zahl wird zur ersten Zahl
  \item[\texttt{b = c;}] \hfill\\
    und die neu berechnete nächste Zahl wird zur zweiten Zahl.
  \item[\texttt{printf("{}Fibonacci \%d: \%d\backslash n"{}, i, a);}] \hfill\\
    Wir geben die aktuelle Fibonacci-Zahl auf dem Bildschirm aus. Das \texttt{\%d} ist ein Platzhalter für die aktuelle Zahl, die an dieser Stelle ausgegeben wird. 
\end{description}

\paragraph{Merke:}

\begin{itemize}[noitemsep]
    \item Die \texttt{for}-Schleife ermöglicht es einem, Code gemäß einer Bedingung beliebig oft zu wiederholen.
\end{itemize}
